% Independent contribution. AI-produced by GPT-6 Astra; no priority claimed.
\subsection{Marked prime-power quotients of a genus-zero surface}

Here a marked quotient retains the image of every ramification
point as an individually named branch point.  The assertion
below concerns the GT parameter of its mapping class action;
geometric permutations of branch names are retained as the
ordinary symmetric-group factor.  It does not assert a result
for a quotient whose covering surface has positive genus.

\begin{result}{B.4.1}{Proved}{Genus-zero prime-power quotients with branch markings}
Let a finite prime-power group act faithfully and preserve
orientation on a punctured sphere.  Mark all its ramification
points and, when needed, adjoin regular orbits.  The action
induced from the ordinary full GT tower on the completed
normalizer quotient of this marked cover has the same full GT
parameter as the original action.  The assertion is compatible
with the ordinary maps forgetting the adjoined regular orbits.
\end{result}
\begin{proof}
Nielsen realization and uniformization realize the action as
M\"obius transformations \cite[Theorem 5 and \S IV]{CA_Kerckhoff1983}.
Here the prime-power classification has a short direct proof.
Choose a central element of order $p$.  Its fixed set consists
of two points, so the whole group preserves this pair.  The
kernel of its permutation action on the pair fixes both points,
and after sending them to $0,\infty$ it is a finite subgroup of
$z\mapsto az$, hence cyclic.  For odd $p$ the permutation action
is trivial.  For $p=2$ its nontrivial coset is represented by
$z\mapsto b/z$, an involution which inverts the cyclic kernel;
rescaling $z$ gives $z\mapsto z^{-1}$.  This proves the cyclic
or dihedral alternatives by exact actions.  The cyclic case,
with the actual quotient and comparison maps, is B.1.7.  We
give the maps for the dihedral case.

Write
\[
 H=\langle r,t\mid r^N=t^2=1,\ trt=r^{-1}\rangle,
 \qquad N=2^a,
 \qquad K=\langle r\rangle.
                                                               \tag{B.L1}
\]
Use the explicit action $r(z)=\zeta_N z$, $t(z)=z^{-1}$.
Its branch orbits are
\[
 \{0,\infty\},\qquad \mu_N,\qquad
              \{z:z^N=-1\}.                              \tag{B.L2}
\]
Adjoin at least two named regular $H$-orbits, so that each of
both intermediate sphere quotients has enough regular orbits
for the cyclic parameter test.  The first quotient is the
actual map
\[
 S\longrightarrow S/K,\qquad u=z^N.                       \tag{B.L3}
\]
Its fixed branch marks are $u=0,\infty$.  The two reflection
branch orbits in (B.L2) give the individual marks $u=1,-1$.
Every regular $H$-orbit gives a pair of marked points
$\{v,v^{-1}\}$ in the quotient.  The residual subgroup
$\overline H=H/K=C_2$ acts by $u\mapsto u^{-1}$, with the
marked fixed points $1,-1$.  Thus its other marked orbits
are exactly the pair $\{0,\infty\}$ and the pairs
$\{v,v^{-1}\}$.  The second quotient can be written
\[
 S/K\longrightarrow S/H,\qquad w=u+u^{-1}.                \tag{B.L4}
\]
In particular the two displayed coverings, their branch sets,
and their composition are specified before completion.

Let $G$ be the relevant mixed mapping class group upstairs.
For $N\geq4$, $K$ is the characteristic rotation subgroup of
$H$, so $N_G(H)\subseteq N_G(K)$.  The standard lifting map
identifies
\[
 G_K=N_G(K)/K
\]
with the finite-index mixed sphere mapping class group which
preserves the two cyclic branch labels as an unordered pair
and the displayed remaining label blocks.  The cover (B.L3)
realizes this map.  It contains the pure sphere mapping class
group on all these labels.  Every pure quotient mapping
class lifts: it preserves the cyclic monodromy, and a branch
interchange acts on $K$ by inversion.  Since $H$ is precisely
the inverse image of $\overline H$ under $N_G(K)\to G_K$, we
have the exact discrete equalities and quotient isomorphism
\[
 \begin{split}
 N_G(H)/K&=N_{G_K}(\overline H),\\
 N_G(H)/H&\xrightarrow{\ \sim\ }
                      N_{G_K}(\overline H)/\overline H.
 \end{split}                                                \tag{B.L5}
\]
For example the reverse inclusion in the first equality
follows by lifting a normalizer of $\overline H$ and taking
the inverse image of that subgroup; no pointwise-centralizer
identity is substituted for this normalizer statement.
The full normalizer topologies proved in Part A.1, together
with the finite kernels in (B.L5), give these same maps on
full profinite completions.

By B.1.7, the restriction of the induced action on $G_K$ to
its pure sphere subgroup has GT parameter $F$.  Its geometric
permutation factor is the one from the retained labels.  The
extension to the mixed group is determined by this restriction:
after the one pure inner adjustment making two restrictions
agree literally, their values at any mixed element differ by
an element centralizing the pure subgroup, and that centralizer
is trivial.  Indeed there are at least eight intermediate marks;
the induced permutation embeds in the symmetric-group factor of
$\operatorname{Out}(\widehat\Gamma_{0,n})$ by HMM, so a
centralizing element has trivial permutation, and the pure
mapping class group has trivial center.  Thus it is the ordinary mixed genus-zero action
of the same parameter, with the displayed geometric relabelling.
Apply the cyclic $N=2$ comparison B.1.7 to the actual residual
action $u\mapsto u^{-1}$ and (B.L4).  Its quotient parameter
is therefore again $F$.  The isomorphism (B.L5) identifies
this with the original dihedral quotient parameter.

When $N=2$, $H$ is the Klein four group.  Its three order-two
subgroups are permuted by its geometric normalizer, whose
quotient action is the full $\mathfrak S_3$ on these subgroups
(the explicit Klein-four normalizer was computed in Part A.1).
After that geometric correction one may retain a chosen
$K=C_2$ in (B.L1).  Its resulting permutation of the three
branch names is kept in the symmetric-group factor.  The
same two quotient maps and inverse-image proof (B.L3)--(B.L5)
then apply to the normalizer which retains $K$.  This subgroup
has finite index in the full normalizer.  Its pure quotient
is the full pure mapping class group of the marked quotient
sphere, and its GT parameter is $F$ by the two cyclic steps.
Consequently the full normalizer has the same GT parameter;
its remaining geometric branch permutations commute with
that factor.  This also agrees with the direct Klein-four
comparison already established in Part A.1.

Finally, forgetting an adjoined regular $H$-orbit commutes
with both power coverings and the map $u\mapsto u+u^{-1}$.
The ordinary tower gives the corresponding homomorphisms of
normalizers.  Making the finite-subgroup correction on the
largest labelled surface once and then forgetting gives
compatible corrections on every smaller surface.  The
action assertion therefore descends through these maps, as claimed.
For a smaller quotient this means equality with the ordinary
$F$-action under the displayed forget map.  Uniqueness or
recovery of a GT parameter is asserted only at the retained
parameter-detecting comparison, such as
$\widehat\Gamma_{0,5}$; the trivial group
$\widehat\Gamma_{0,3}$ does not recover a parameter.
\end{proof}
